\chapter{Fonctions affines et comparaisons}
\section*{Exercice 13.1 -- Reconstruire}
Une fonction affine vérifie $f(1)=5$ et $f(4)=14$. Déterminer son expression.

\section*{Exercice 13.2 -- Signe et variation}
Étudier le signe et les variations de $f(x)=-4x+12$.

\section*{Exercice 13.3 -- Comparer des tarifs}
Pour $x\geq0$, comparer $A(x)=30+2x$ et $B(x)=10+6x$.

\section*{Exercice 13.4 -- Une preuve}
Prouver qu'une fonction affine de pente positive est strictement croissante.

\section*{Exercice 13.5 -- Interpolation}
On suppose une évolution affine de $18$ à $30$ entre les instants $t=2$ et $t=8$. Estimer la valeur à $t=5$. Préciser l'hypothèse.

\section*{Exercice 13.6 -- Modèle physique}
Un réservoir de $75$ litres perd $2{,}5$ litres par minute. Donner la fonction de volume, son domaine utile et l'instant où il reste $20$ litres.
